% !TEX root = ../../../main.tex
% Appendix entry: catalogue.gramian_gasf -- "Gramian GASF" (Appendix A).
\entryfield{Algorithm}
The block wraps \code{compute_GASF} in \code{Working/Catalogue/gramian/gramian_calc.py}, the
Gramian angular summation field of \citet{wang2015imaging}. For a span $x_0,\dots,x_{n-1}$:
\begin{enumerate}
\item \emph{Rescale} (\code{_polar_encode}) to $[-1, 1]$ by the span's own extremes:
  \[ \tilde x_i = \frac{2\,(x_i - \min_k x_k)}{\max_k x_k - \min_k x_k} - 1, \]
  clipped to $[-1,1]$ against rounding.
\item \emph{Angle}: $\phi_i = \arccos \tilde x_i \in [0, \pi]$.
\item \emph{Field}:
  \[ G_{ij} = \cos(\phi_i + \phi_j) = \tilde x_i \tilde x_j - \sqrt{1-\tilde x_i^2}\,\sqrt{1-\tilde x_j^2}, \]
  computed as \code{np.cos(np.add.outer(phi, phi))}.
\end{enumerate}
The paper's polar radius (the time stamp) plays no part in the field and is not computed.
There is \emph{no piecewise aggregate approximation (PAA)}: \citet{wang2015imaging} shrink the
series with PAA before building the image, but this code builds the full $n \times n$ field
from every sample. Any shrinking to a network's input size happens later, as image resizing,
in \blockref{catalogue.window_images}.

\entryfield{Assumptions}
\begin{itemize}
\item \emph{Units.} Not unit-aware, and it does not need to be: the min--max rescaling
  removes level and scale, so a span in volts and the same span in millivolts give the same
  image.
\item \emph{Sampling rate.} Ignored. One pixel is one sample, whatever the rate.
\item \emph{Missing values.} \code{x.min()} and \code{x.max()} are NaN if any sample is NaN,
  so a single missing sample makes every pixel NaN.
\item \emph{Span length.} At most \code{max_span_samples} $= 5000$ (below).
\end{itemize}

\entryfield{Outputs}
An \code{Encoding} with \code{kind="image"} whose values are an $n \times n$ float64 array in
$[-1, 1]$, symmetric ($G_{ij} = G_{ji}$), with the diagonal $\cos 2\phi_i = 2\tilde x_i^2 - 1$.
The adapter adds nothing to \code{meta}. The block's plot hook (\code{plot_gramian_suite})
draws all four encodings of the span together, not only this one. Like any step, the result
is kept by the executor's step cache when it is slow enough to be worth it
(\code{Working/execution.py}).

\entryfield{Cost}
$O(n^2)$ time and memory: the field is an $n \times n$ float64 matrix, 8 bytes a cell. The
adapter declares \code{max_span_samples=5000} (\code{MAX_SPAN_SAMPLES} in
\code{Adapters/catalogue_gramian_gasf.py}), which caps the matrix at $5000^2 \times 8$ bytes
$= 200$\,MB. \code{execute_recipe} in \code{Working/execution.py} checks the span against it
before calling the block and refuses a longer span with a message naming the limit. The
same number is duplicated as \code{WM_GRAMIAN_MAX_SAMPLES} in \code{Working/config.py} for the
window matrix's image stage, and a test keeps the two equal.

\entryfield{Failure modes}
\begin{itemize}
\item \emph{A flat span} (all samples equal) divides zero by zero: every pixel is NaN.
\item \emph{Any NaN in the span}: every pixel is NaN (above).
\item \emph{One outlier sets the scale.} Because the rescaling uses the span's extremes, one
  large spike pushes every other sample towards one end of $[-1,1]$, and the rest of the
  image loses contrast.
\item \emph{Different spans are not on one scale.} Two spans with very different amplitudes
  can give the same image, because each is rescaled by its own extremes.
\end{itemize}

\entryfield{Validation}
The formula is the published one \citep{wang2015imaging}. No test in the repository checks
the values of the image: the tests check the block's declared types and that a too-long span
is refused, and one checks the shape of the result ($600 \times 600$ for a 600-sample span).

\entryfield{References}
\citet{wang2015imaging}.

\entryfield{Tests}
\code{tests/test_t07_adapter_remap.py} checks the declared input and output types and that
\code{execute_recipe} refuses a span one sample over the limit.
\code{tests/test_window_matrix_store.py} checks that the window matrix's limit equals this
block's, refuses a whole-channel span and checks the output shape on a 600-sample span (the
last two skip when the real channel data are absent). \code{tests/test_execution.py},
\code{tests/test_chain_validation.py} and \code{tests/test_artifacts.py} use the block as a
stand-in step in a recipe.
